\chapter{The project skeleton}
\label{ch:project}

\chapterbrief
  {Copy the template, build it, run its tests --- all on your laptop.}
  {What Maven's folder layout means, and why it is not negotiable.}
  {\code{Tests run: 3, Failures: 0}.}

\section{Copy the template}

\begin{shell}
# on your laptop
cp -r starter-pack/template ~/midnight-brewer
cd ~/midnight-brewer
ls -la
\end{shell}

\section{Build it before you change anything}

\begin{shell}
mvn test
\end{shell}

The first run downloads dependencies and takes a minute or two. Afterwards it is
seconds.

\begin{trace}
[INFO] Tests run: 3, Failures: 0, Errors: 0, Skipped: 0
[INFO] BUILD SUCCESS
\end{trace}

\begin{trap}{Read the number, not the word}
\code{BUILD SUCCESS} with \emph{Tests run: 0} is also a success, and it means
your tests are not being found at all.

This is the single most common silent failure in Java projects: a misplaced test
folder, a class not ending in \code{Test}, an old Surefire version that cannot
see JUnit 5. Everything looks green and nothing is verified.

\textbf{Always check the count.} Here it must say 3.
\end{trap}

\section{What is in the folder}

\begin{center}
\begin{tikzpicture}[font=\footnotesize\ttfamily, node distance=0mm,
  f/.style={anchor=west, inner sep=1.5pt},
  c/.style={anchor=west, inner sep=1.5pt, font=\footnotesize\sffamily\itshape,
            color=brewergrey}]
  \def\y{0}
  \node[f] at (0,0.0)   {midnight-brewer/};
  \node[f] at (0.4,-0.5) {pom.xml};
    \node[c] at (4.6,-0.5) {what to build, which libraries};
  \node[f] at (0.4,-1.0) {Makefile};
    \node[c] at (4.6,-1.0) {the commands you type};
  \node[f] at (0.4,-1.5) {.gitignore};
    \node[c] at (4.6,-1.5) {keeps target/ out of git};
  \node[f] at (0.4,-2.0) {README.md};
    \node[c] at (4.6,-2.0) {how to use it};
  \node[f] at (0.4,-2.5) {scripts/};
  \node[f] at (0.9,-3.0) {pi-bootstrap.sh};
    \node[c] at (4.6,-3.0) {run once, on the Pi};
  \node[f] at (0.9,-3.5) {deploy.sh};
    \node[c] at (4.6,-3.5) {build here, copy there};
  \node[f] at (0.9,-4.0) {run.sh};
    \node[c] at (4.6,-4.0) {run a class on the Pi};
  \node[f] at (0.4,-4.5) {src/main/java/...};
    \node[c] at (4.6,-4.5) {your program};
  \node[f] at (0.4,-5.0) {src/test/java/...};
    \node[c] at (4.6,-5.0) {your tests};
  \node[f] at (0.4,-5.5) {target/};
    \node[c] at (4.6,-5.5) {build output -- never edited, never committed};
\end{tikzpicture}
\end{center}

\begin{concept}{Why \texttt{src/main/java} and not something shorter}
Maven follows \emph{convention over configuration}: it knows where sources live,
so you never write it down.

Put a file in the wrong place and it is silently not compiled --- no error, it
simply does not exist as far as the build is concerned. That is the trade: total
consistency across every Maven project in the world, at the cost of paths being
fixed.

The package \code{com.midnightbrewer.diag} must map to the folder
\code{src/main/java/com/midnightbrewer/diag/}. Not similar --- identical.
\end{concept}

\section{Run the toolchain check}

\begin{shell}
mvn -q compile
java -cp target/classes com.midnightbrewer.ToolchainCheck
\end{shell}

\begin{trace}
=== toolchain check ===
java.version    = 21.0.11
os.name         = Mac OS X
os.arch         = aarch64
major version   = 21
sanity          = 0x92 (146)

  -> toolchain OK. Nothing above needed hardware.
\end{trace}

\begin{concept}{What \texttt{-cp} means}
\code{-cp} is the \emph{classpath}: where Java looks for compiled classes.

\code{target/classes} is where Maven put your compiled code. Later you will use
\code{-cp "target/classes:target/dependency/*"} to include the library jars as
well --- the colon separates entries, and the \code{*} means ``every jar in that
folder''.

Most ``\code{ClassNotFoundException}'' and ``\code{NoClassDefFoundError}''
problems are a classpath missing an entry, not a missing file.
\end{concept}

\section{Read the test}

Open \code{src/test/java/com/midnightbrewer/ToolchainCheckTest.java}. One test
is worth understanding now, because it encodes the mistake you are most likely
to make in Volume One:

\begin{jcode}
@Test
void signedByteMaskingWorks() {
    byte b = (byte) 0x92;

    assertEquals(-110, b, "a signed byte holding 0x92 reads as -110");
    assertEquals(146, b & 0xFF, "masked with 0xFF it reads as 146");
}
\end{jcode}

Java's \code{byte} is signed and there is no \code{unsigned} keyword. \hex{0x92}
has its top bit set, so read as a number it is $-110$. The idiom \code{\& 0xFF}
promotes to \code{int} and strips the sign extension.

You will forget this at least once. Everyone does.

\section{The Makefile}

\begin{shell}
make            # lists every target
\end{shell}

\begin{trace}
Midnight Brewer -- available targets:

  build        compile locally
  test         run the tests locally
  deploy       build + copy classes to the Pi
  run          run MAIN on the Pi
  probe        run the RC522 wiring probe on the Pi
  pi-setup     copy + run the one-time Pi bootstrap script
  ssh          open a shell on the Pi
\end{trace}

\begin{concept}{Why a Makefile and not just typing the commands}
Because you will type these hundreds of times, and because six months from now
\code{make deploy} still works when you have forgotten the rsync flags.

A Makefile is executable documentation. The commands are recorded in the
repository, they are the same for everyone, and \code{make} with no arguments
tells you what is available --- which is more than most projects manage.
\end{concept}

\section{Point it at your Pi}

\begin{shell}
export PI=martin@brewer
\end{shell}

Add that to \code{\textasciitilde/.zshrc} or \code{\textasciitilde/.bashrc} so it
survives a new terminal. It uses the \code{brewer} alias from your
\code{\textasciitilde/.ssh/config}.

\section{Commit it}

Now do Chapter~\ref{ch:git} if you have not:

\begin{shell}
git init
git add .
git commit -m "chore: project skeleton from the starter pack"
\end{shell}

\begin{checkpoint}{The project builds}
\begin{shell}
cd ~/midnight-brewer
mvn test
make
\end{shell}

\code{Tests run: 3, Failures: 0}, and \code{make} lists its targets.

No hardware has been involved yet. That is the point: when something breaks in
Volume One, you already know the toolchain is not the problem.
\end{checkpoint}
